\cgeentry{cge_balance}{Balance a SAM by RAS or cross entropy}

\cgesec{Description}

\cmd{cge_balance} makes the SAM of the current model balanced, so that every
account's row total equals its column total, and starts a new model from the
result; any earlier mapping, calibration and simulations are discarded. Two
methods are available: RAS, which sets every account's total to a target, and
cross entropy, which estimates the totals that are not known together with the
cells. Both keep the sign of every cell, every zero cell and the cells listed
in \opt{fix}{ed()}, and both minimise the same measure of the information added
to the SAM.

\cgesec{Quick start}

Balance the SAM by RAS, with each account's target the average of its row and
column totals
\begin{stcode}
cge_balance
\end{stcode}
As above, with the row totals as targets and two totals given
\begin{stcode}
cge_balance, targets(rows) totals(act_agr:97 hhd:290)
\end{stcode}
Balance by cross entropy, with the household's total known, keeping every
payment to the rest of the world and every import tariff
\begin{stcode}
cge_balance, method(ce) totals(hhd:290) fixed(row:* tax_imp:*)
\end{stcode}

\cgesec{Syntax}

\begin{stsyntax}
\cmdb{cge_balance} [\texttt{,} \argu{options}]
\end{stsyntax}

\begin{opttable}
\optgroup{Main}
\optarg{meth}{od}{method} & \cmd{ras} or \cmd{ce}; default \cmd{ras}\\
\optarg{targ}{ets}{rule} & RAS targets of the accounts not in \opt{tot}{als()}: \cmd{average}, \cmd{rows} or \cmd{columns}; default \cmd{average}\\
\optarg{tot}{als}{\argu{account}\texttt{:}\argu{value} \dots} & known account totals\\
\optarg{fix}{ed}{list} & cells kept at their values\\
\optgroup{Iteration}
\optarg{tol}{erance}{\#} & tolerance, relative to the largest absolute cell; default \texttt{1e-10}\\
\optarg{iter}{ate}{\#} & maximum number of iterations; default 10,000 for RAS and 100 for cross entropy\\
\optgroup{Reporting}
\opt{nod}{isplay} & suppress the table of the largest changes\\
\end{opttable}

\cgesec{Options}

\begin{optdesc}
\item[\optarg{meth}{od}{method}] chooses the method: \cmd{ras}, generalised
RAS (the default), or \cmd{ce}, cross entropy.

\item[\optarg{targ}{ets}{rule}] sets the RAS targets of the accounts not given
in \opt{tot}{als()}: \cmd{average}, the average of the account's row and column
totals (the default); \cmd{rows}, its row total; or \cmd{columns}, its column
total. It applies only to RAS: with \cmd{method(ce)}, \cmd{rows} and
\cmd{columns} are errors, and known totals are given in \opt{tot}{als()}.

\item[\optarg{tot}{als}{\argu{account}\texttt{:}\argu{value} \dots}] gives
known account totals, separated by spaces: targets for RAS, constraints for
cross entropy. For example, \cmd{totals(act_agr:97 hhd:290)} gives the totals
of two accounts. An account may be given only once.

\item[\optarg{fix}{ed}{list}] lists, separated by spaces, the cells kept at
their values: a cell \argu{row}\cmd{:}\argu{col}, a whole row
\argu{row}\cmd{:*}, a whole column \cmd{*:}\argu{col}, or an account's row and
column, \argu{account}. Every account named must be in the SAM; \cmd{*:*} is an
error.

\item[\optarg{tol}{erance}{\#}] sets the convergence tolerance: every balance
and every target or known total is met to \argu{\#} times the larger of one
and the largest absolute cell of the SAM before balancing. The default is
\texttt{1e-10}, well within the balance that \manref{cge_map} requires.

\item[\optarg{iter}{ate}{\#}] sets the maximum number of iterations: RAS
iterations, each a row and a column step, or Newton steps of cross entropy. The
default is 10,000 for RAS and 100 for cross entropy.

\item[\opt{nod}{isplay}] suppresses the table of the cells with the largest
changes.
\end{optdesc}

\cgesec{Remarks and examples}

SAMs assembled from national accounts, supply--use tables, household surveys
and government budgets rarely balance. \cmd{cge_balance} adjusts the cells so
that every account balances, changing the SAM as little as possible in the
sense of the measure given in Methods and formulas. Neither method changes the
sign of a cell or turns a zero into a flow, so structural zeros stay zero, and
cells in \opt{fix}{ed()} keep their values. The methods differ in what is known
about the totals:
\begin{itemize}
\item RAS needs a target for every account's total. The default target is the
average of the account's row and column totals; \opt{targ}{ets()} and
\opt{tot}{als()} choose others. Negative cells are handled by generalised RAS
(GRAS).
\item Cross entropy needs no targets. It estimates every total that is not
given in \opt{tot}{als()} together with the cells, subject to balance. With
every total given, it gives the RAS solution.
\end{itemize}
\cmd{cge_balance} displays the number of iterations, the largest imbalance and
the SAM total before and after, the number of cells changed and the largest
relative change of a cell, and the ten cells with the largest absolute changes,
followed by the size, total and imbalance of the new SAM, as \manref{cge_sam}
displays them.

The balanced SAM is an estimate. Which totals were given, which cells were
fixed and which method was used should be stated with any results from it.

\cgesubsec{Example 1: RAS}

The unbalanced synthetic SAM is out of balance by up to 12. RAS to the average
totals balances it:
\stlog{bal_ras}
The largest imbalance after balancing, \texttt{1.53e-08}, is within the
tolerance: \texttt{1e-10} times the largest cell, 158.729. With the average
targets, the SAM total is unchanged, since the targets add up to it. Every
non-zero cell may change, and every one does.

\cgesubsec{Example 2: cross entropy with known totals and fixed cells}

Suppose that the output of agriculture, 97, and household income, 290, are
known, and that the payments to the rest of the world and households' labour
income are taken as reliable. Cross entropy estimates every other total:
\stlog{bal_ce}
The eight fixed cells, the seven non-zero payments to the rest of the world
and \cmd{hhd:lab}, keep their values, and so do the transfers of the enterprise
and of the government to themselves: a diagonal cell is in both the row and
the column of its account, and balance alone does not change it. The other 58
non-zero cells change. Household income of 290 is above both the household's
row total, 247.9, and its column total, 255.1, in the unbalanced SAM, so its
receipts and its payments rise, and the SAM total with them.

The results are stored in \cmd{r()}:
\stlog{bal_return}

\cgesubsec{Example 3: a SAM that cannot be balanced}

Zeros and fixed cells may make the targets impossible to meet. Fixing the
household's row and column leaves its two totals as they are, and they differ:
\stlog{bal_fixed}
\cmd{cge_balance} checks such contradictions before it starts: an account whose
row and column cannot change but differ, or a given total that differs from
the total of cells that cannot change. Other contradictions arise from the
pattern of zeros and fixed cells. For example, in a SAM whose activities and
product-tax accounts are paid only by commodities, fixing every import cell
fixes the difference between the commodities' totals and theirs; RAS targets
that do not respect it cannot be met, and RAS stops with error 430 when its gap
no longer falls. Cross entropy, which needs fewer totals, can often balance
such a SAM.

\cgesec{Stored results}

\cmd{cge_balance} stores the following in \cmd{r()}:
\begin{storedres}
\resgroup{Scalars}
\res{r(iterations)} & RAS iterations or Newton steps of cross entropy\\
\res{r(maxdiff_before)} & largest row--column imbalance before balancing\\
\res{r(maxdiff)} & largest row--column imbalance after balancing\\
\res{r(total_before)} & sum of all cells before balancing\\
\res{r(total)} & sum of all cells after balancing\\
\res{r(N_changed)} & number of cells changed\\
\res{r(maxrelchange)} & largest relative change of a cell\\
\res{r(N_accounts)} & number of accounts\\
\res{r(N_blank)} & always 0\\
\resgroup{Macros}
\res{r(method)} & \cmd{ras} or \cmd{ce}\\
\end{storedres}

\cgesec{Methods and formulas}

Let $\bar X$ be the SAM before balancing and $X$ the SAM after it. A cell may
change if it is not zero and not in \opt{fix}{ed()}; let $Z$ be the set of such
cells, and write $P$ for the positive cells of $\bar X$ in $Z$, $N$ for the
absolute values of its negative cells in $Z$, and $F$ for all other cells,
which keep their values. Both methods minimise
\begin{equation*}
D(X, \bar X) = \sum_{(i,j) \in Z} \Bigl( |x_{ij}| \ln \frac{|x_{ij}|}{|\bar x_{ij}|} - |x_{ij}| + |\bar x_{ij}| \Bigr),
\end{equation*}
a measure of the information added to $\bar X$ (Junius and Oosterhaven 2003;
Lenzen, Wood and Gallego 2007), which for a SAM without negative cells is the
cross entropy that RAS minimises (McDougall 1999). It is not the
column-coefficient cross entropy of Robinson, Cattaneo and El-Said (2001),
which measures the distance between the column shares of $X$ and $\bar X$;
that method is not implemented. Convergence is measured against
$\sigma = \max(1, \max_{ij} |\bar x_{ij}|)$: the iterations stop when every
balance, target and known total is met to $\text{\opt{tol}{erance()}} \times
\sigma$.

\textit{RAS.} Every account $k$ has a target total $t_k$: the average of its
row and column totals in $\bar X$, one of them as \opt{targ}{ets()} says, or
the value in \opt{tot}{als()}. If an account's row has no cell that may change,
its target is its row total, and likewise for its column; if both are fixed and
differ, or a given total contradicts them, there is no balanced SAM and
\cmd{cge_balance} stops. With $f^r_i = \sum_j f_{ij}$ and $f^c_j = \sum_i
f_{ij}$, the row and column totals of the cells that may change must be $u_i =
t_i - f^r_i$ and $v_j = t_j - f^c_j$. Generalised RAS finds row factors $r_i$
and column factors $s_j$ such that
\begin{equation*}
x_{ij} = r_i\, p_{ij}\, s_j - \frac{n_{ij}}{r_i\, s_j} + f_{ij}
\end{equation*}
has these totals. Starting from $s = 1$, it alternates a row step and a column
step. The row step solves $r_i p_i - n_i / r_i = u_i$ for each row, with
$p_i = \sum_j p_{ij} s_j$ and $n_i = \sum_j n_{ij} / s_j$:
\begin{equation*}
r_i = \frac{u_i + \sqrt{u_i^2 + 4 p_i n_i}}{2 p_i} \quad (u_i \ge 0), \qquad
r_i = \frac{2 n_i}{\sqrt{u_i^2 + 4 p_i n_i} - u_i} \quad (u_i < 0),
\end{equation*}
two forms of the same root that avoid cancellation. The column step is the
same with the roles of rows and columns exchanged. Without negative cells
this is RAS. A factor that is not positive and finite means that the target
cannot be met with the signs of the cells, and \cmd{cge_balance} stops. After
100 iterations, RAS also stops, with error 430, when the gap has not fallen by
a relative $10^{-6}$ over the last 100 iterations.

\textit{Cross entropy.} Let $K$ be the accounts in \opt{tot}{als()}. $X$
minimises $D(X, \bar X)$ subject to
\begin{equation*}
\sum_j x_{kj} - \sum_i x_{ik} = 0 \quad \text{for every account } k, \qquad
\sum_j x_{kj} = t_k \quad \text{for } k \in K.
\end{equation*}
The solution has the form
\begin{equation*}
x_{ij} = p_{ij}\, e^{a_{ij}} - n_{ij}\, e^{-a_{ij}} + f_{ij}, \qquad
a_{ij} = \mu_i - \mu_j + \rho_i,
\end{equation*}
with a multiplier $\mu_k$ for the balance of each account and $\rho_k$ for
each known total ($\rho_k = 0$ for $k \notin K$). The multipliers minimise the
convex dual
\begin{equation*}
\varphi(\mu, \rho) = \sum_{ij} p_{ij}\, e^{a_{ij}} + \sum_{ij} n_{ij}\, e^{-a_{ij}}
+ \sum_k \mu_k \bigl(f^r_k - f^c_k\bigr) + \sum_{k \in K} \rho_k \bigl(f^r_k - t_k\bigr),
\end{equation*}
whose gradient is the vector of the balances and of the differences from the
known totals, and whose Hessian is
\begin{equation*}
H = \sum_{ij} \bigl(p_{ij}\, e^{a_{ij}} + n_{ij}\, e^{-a_{ij}}\bigr)\, d_{ij} d_{ij}', \qquad
d_{ij} = \frac{\partial a_{ij}}{\partial (\mu, \rho)}.
\end{equation*}
Newton's method minimises $\varphi$, with a backtracking line search that
halves the step until $\varphi$ falls by at least $10^{-4}$ times the decrease
the step predicts. Adding a constant to every $\mu_k$ in a set of accounts
connected by cells that may change leaves $X$ unchanged, so one $\mu$ in each
such set is fixed at zero; the fixed cells of each set must balance among
themselves. A known total whose row or column cannot change must equal the
total of its fixed cells, and needs no multiplier. If Newton's method does not
converge, no balanced SAM with the same pattern of signs and zeros may exist.
When every total is known, the solution is the RAS solution, with $r_i =
e^{\mu_i + \rho_i}$ and $s_j = e^{-\mu_j}$.

\textit{Results.} The number of cells changed counts the cells with $x_{ij} \ne
\bar x_{ij}$, and the largest relative change is $\max_{ij} |x_{ij} - \bar
x_{ij}| / |\bar x_{ij}|$ over the non-zero cells of $\bar X$.
Balancing is described in \secref{sec:balancing}, RAS in \secref{sec:ras} and
cross entropy in \secref{sec:ce}.

\cgesec{References}

\begin{list}{}{\setlength{\leftmargin}{2em}\setlength{\itemindent}{-2em}%
  \setlength{\itemsep}{2pt}\setlength{\parsep}{0pt}}
\item Junius, T., and J. Oosterhaven. 2003. The solution of updating or
regionalizing a matrix with both positive and negative entries.
\textit{Economic Systems Research} 15: 87--96.
\item Lenzen, M., R. Wood, and B. Gallego. 2007. Some comments on the GRAS
method. \textit{Economic Systems Research} 19: 461--465.
\item McDougall, R. A. 1999. Entropy theory and RAS are friends. GTAP Working
Paper 6, Center for Global Trade Analysis, Purdue University.
\item Robinson, S., A. Cattaneo, and M. El-Said. 2001. Updating and estimating
a social accounting matrix using cross entropy methods. \textit{Economic
Systems Research} 13: 47--64.
\end{list}

\cgesec{Also see}

\begin{alsosee}
\seeentry{cge_sam}{Read a social accounting matrix}
\seeentry{cge_aggregate}{Merge SAM accounts into groups}
\seeentry{cge_samexport}{Write the SAM to a file or to memory}
\seeentry{cge_map}{Assign SAM accounts to model roles}
\seechapter{chap:samtools}{SAM tools}
\end{alsosee}
